%% ========================================================================
%% Appendix E: Development Environment Installation Guide
%% ========================================================================

\chapter{Development Environment Installation Guide}
\label{app:instalasi}

Chapter~1 assumes four command-line tools are already installed: PHP,
Composer, Node.js, and Git (see the tip box right before
\cmd{composer create-project} in that chapter). This appendix walks
through installing all four from scratch, split by operating system, then
closes with a verification step common to all three.

\begin{table}[htbp]
\centering
\caption{Tools required before starting Chapter~1}
\label{tab:instalasi-kebutuhan}
\begin{tblr}{
  colspec = {X[0.8,l] X[0.9,l] X[1.6,l]},
  hline{1,2,Z} = {solid},
}
\textbf{Tool} & \textbf{Minimum Version} & \textbf{Purpose} \\
PHP & 8.2 & Runs Laravel 13 and the Artisan command line \\
Composer & 2.x & Installs PHP dependencies from \file{composer.json} \\
Node.js & 18 LTS & Runs \cmd{npm} for the Vite + Tailwind toolchain (Chapter~3) \\
Git & 2.x & Tracks code versions and follows the \file{chapter-NN} branches \\
\end{tblr}
\end{table}

\section{Windows}

The quickest path on Windows 10/11 is \pkg{winget}, the package manager
already bundled with most modern installations. Open PowerShell and run
the following four commands one at a time:

\begin{lstlisting}[language=bash]
winget install PHP.PHP.8.3
winget install Composer.Composer
winget install OpenJS.NodeJS.LTS
winget install Git.Git
\end{lstlisting}

\begin{warningbox}
The default \file{php.ini} shipped with the \pkg{PHP.PHP.8.3} package
disables several extensions, including \texttt{pdo\_sqlite}, which Simple
POS requires. Run \cmd{php --ini} to find the active \file{php.ini}
path, open it, and remove the leading \texttt{;} from the
\texttt{extension=pdo\_sqlite} and \texttt{extension=fileinfo} lines. See
\cref{sec:instalasi-ekstensi} for the full list of required extensions.
\end{warningbox}

\begin{tipbox}
Close and reopen the terminal window after \pkg{winget} finishes
installing. A newly added PATH entry is only picked up by terminals
opened after the change, not the one still running.
\end{tipbox}

Another option is \pkg{Scoop}, a package manager that installs every tool
into the user's own directory (\file{\textasciitilde\textbackslash scoop})
without requiring administrator rights or a UAC confirmation prompt. This
is useful on campus lab or office computers whose accounts are not granted
administrator access. Open PowerShell and run:

\begin{lstlisting}[language=bash]
Set-ExecutionPolicy -ExecutionPolicy RemoteSigned -Scope CurrentUser
irm get.scoop.sh | iex
scoop install php composer nodejs-lts git
\end{lstlisting}

\begin{notebox}
The first two lines install Scoop itself (allowing locally signed
PowerShell scripts to run, then downloading and running Scoop's
installer); the third line installs all four tools at once from Scoop's
default \texttt{main} bucket. Scoop's \pkg{php} package uses the same
default Windows \file{php.ini} as the \pkg{winget} package, so the
\texttt{pdo\_sqlite} extension warning above applies here exactly the
same way. Pick one package manager only, \pkg{winget} or \pkg{Scoop}, not
both; two competing PHP installations on PATH are the most common cause
of \cmd{php -v} showing an unexpected version (see the warning box in the
Verification section).
\end{notebox}

\section{macOS}

\pkg{Homebrew} is the most common third-party package manager on macOS.
If it is not installed yet, follow the instructions at
\url{https://brew.sh} first, then run:

\begin{lstlisting}[language=bash]
brew install php composer node git
\end{lstlisting}

\begin{notebox}
The Homebrew \pkg{php} formula already ships every extension listed in
\cref{sec:instalasi-ekstensi} enabled by default, including
\texttt{pdo\_sqlite}, so the \file{php.ini} edit described for Windows is
not needed on macOS.
\end{notebox}

\begin{warningbox}
macOS ships a system PHP at \file{/usr/bin/php}, usually an old version
Laravel 13 does not support. After \cmd{brew install php}, confirm
\cmd{which php} points to \file{/opt/homebrew/bin/php} (Apple Silicon) or
\file{/usr/local/bin/php} (Intel), not \file{/usr/bin/php}. If it still
points to the system version, add
\texttt{export PATH="/opt/homebrew/bin:\$PATH"} to
\file{\textasciitilde/.zshrc}, then open a new terminal window.
\end{warningbox}

\section{Linux (Debian/Ubuntu)}

\begin{lstlisting}[language=bash]
sudo apt update
sudo apt install php php-cli php-sqlite3 php-mbstring php-xml php-curl unzip
curl -sS https://getcomposer.org/installer | php
sudo mv composer.phar /usr/local/bin/composer
sudo apt install nodejs npm git
\end{lstlisting}

\begin{notebox}
Ubuntu LTS's default package repository sometimes lags a few PHP versions
behind the latest release. If \cmd{php -v} reports a version below 8.2
after \cmd{apt install}, add the \pkg{ondrej/php} PPA
(\cmd{sudo add-apt-repository ppa:ondrej/php}) before re-running
\cmd{apt update \&\& apt install php8.3}. Fedora and Arch Linux users can
use the equivalent package names with \cmd{dnf install php composer
nodejs npm git} or \cmd{pacman -S php composer nodejs npm git}.
\end{notebox}

\section{Verification}

Once one of the paths above is done, run the following four commands in
the terminal to confirm everything is installed and reachable through
PATH:

\begin{lstlisting}[language=bash]
php -v
composer -V
node -v
git --version
\end{lstlisting}

\begin{figure}[htbp]
\centering
\includegraphics[width=0.85\linewidth]{verifikasi-instalasi}
\caption{Sample terminal output after running the four verification
commands in sequence; the exact version numbers on your screen will
differ depending on when each tool was installed, what matters is the
absence of any \texttt{command not found} message}
\label{fig:verifikasi-instalasi}
\end{figure}

\begin{warningbox}
If \cmd{php -v} shows an old version even though a newer one was just
installed, another PHP installation is likely found first on PATH
(common on macOS with the system PHP, or on Windows with a leftover
XAMPP install). Run \cmd{which php} (macOS/Linux) or \cmd{where.exe php}
(Windows) to see the order of locations checked, then adjust PATH so the
correct installation appears first.
\end{warningbox}

\section{PHP Extensions Laravel Requires}
\label{sec:instalasi-ekstensi}

Laravel 13 checks for nine PHP extensions when \cmd{composer
create-project} and \cmd{php artisan serve} run. Most are already
enabled by default on modern PHP installations, but minimal builds
(especially manual builds from source, or a bare \texttt{php:cli} Docker
image) often miss some of them.

\begin{table}[htbp]
\centering
\caption{PHP extensions required for Laravel 13 and Simple POS}
\label{tab:instalasi-ekstensi}
\begin{tblr}{
  colspec = {X[1,l] X[1.8,l]},
  hline{1,2,Z} = {solid},
}
\textbf{Extension} & \textbf{Purpose} \\
OpenSSL & Encrypts sessions, cookies, and the application key \\
PDO \& \texttt{pdo\_sqlite} & Eloquent's connection to the \file{database.sqlite} file \\
Mbstring & UTF-8-safe string handling (product names, transaction labels) \\
Tokenizer & Composer's and Blade's internal parser \\
XML & Parses Composer package configuration \\
Ctype \& JSON & Input validation and API responses (Chapter~9) \\
BCMath & High-precision decimal math for transaction totals \\
Fileinfo & File-type detection during CSV import (Chapter~8) \\
\end{tblr}
\end{table}

\begin{tipbox}
The \cmd{php -m} command prints every extension active on the current
PHP installation. Pair it with \cmd{grep} to check one at a time, for
example \cmd{php -m | grep -i sqlite} to confirm \texttt{pdo\_sqlite} is
active before running \artisan{php artisan migrate:fresh --seed} for the
first time.
\end{tipbox}
